\chapter{Introduction}
\label{sec:introduction}

Field-programmable gate arrays (FPGAs) are increasingly used as datacenter accelerators because they provide programmable hardware with high performance and energy efficiency. In cloud and serverless settings, this programmability enables users to submit hardware applications that are deployed on shared FPGA infrastructure. Partial reconfiguration (PR) makes this model practical by allowing user logic to be loaded into a reconfigurable region while the trusted FPGA shell remains active.

This deployment model creates a security-critical boundary. The user-provided partial bitstream becomes an input to the platform, and loading it into the FPGA fabric may expose shared physical resources such as power delivery, thermal headroom, routing infrastructure, and neighboring workloads. Malicious circuits based on ring oscillators (ROs) are a representative threat: at scale, RO arrays can generate high-frequency switching activity, waste power, induce voltage drops, and cause denial-of-service or fault effects.

Existing cloud FPGA platforms show that bitstream validation is a practical concern. Prior attacks on AWS F1 demonstrated that malicious power-hammering designs could pass platform checks and crash instances. More recent work on AWS F2 suggests that HLS-generated power-wasting designs can still bypass deployment flows. These examples motivate inspection mechanisms that complement static platform checks and operate before partial reconfiguration.

Machine-learning-based bitstream inspection is a promising direction because prior work has shown that bitstream-derived representations contain learnable structure and can be used to classify malicious designs. However, most ML-based approaches are evaluated as offline software classifiers. This leaves a systems gap: a serverless FPGA platform needs a checker that fits the deployment path. Such a checker should be accurate enough to reject suspicious bitstreams, small enough to coexist with platform and user logic, and fast enough to avoid dominating PR latency.

This project studies whether ML-based malicious-bitstream inspection can move into the FPGA deployment path itself. Instead of relying on host-side CPU or GPU inspection, we explore a compact FPGA-resident checker for Coyote-style serverless FPGA systems. The checker receives a user partial bitstream, applies a deterministic byte-stream preprocessing step, runs a synthesized CNN classifier, and produces an accept/reject decision or suspiciousness score before PR.

We evaluate this idea using a dataset of benign Coyote application partial bitstreams and RO-based malicious partial bitstreams. Each bitstream is transformed into a fixed-size image-like tensor through deterministic resizing. We then train CNN classifiers, explore input resolution and model depth, synthesize selected candidates with hls4ml, and evaluate the resulting hardware designs across detection quality, false-positive and false-negative behavior, latency, and FPGA resource overhead.

This report makes four contributions:

\begin{enumerate}
    \item \textbf{Partial-bitstream dataset.}
    We construct a Coyote-oriented dataset from benign application partial bitstreams and RO-based malicious payloads.

    \item \textbf{Datacenter FPGA setting.}
    We evaluate inspection in a Coyote-style deployment setting with partial reconfiguration, multitenancy, and UltraScale+-class FPGA hardware.

    \item \textbf{ML-to-hardware inspection path.}
    We connect CNN training to hls4ml synthesis and evaluate whether trained bitstream classifiers can become practical FPGA-resident checkers.

    \item \textbf{Hardware-aware design-space evaluation.}
    We compare candidate models across detection quality, FNR/FPR, ROC AUC, LUT usage, BRAM usage, DSP usage, and added inspection latency.
\end{enumerate}

The main result is that FPGA-resident ML bitstream inspection is feasible as a deployment-path defense-in-depth mechanism. Among the synthesized candidates, the 256$\times$7 CNN provides the preferred deployment point because it balances useful detection quality with lower latency and substantially lower BRAM overhead than the larger 512$\times$7 alternative. Overall, the prototype shows a practical path from user-submitted partial bitstreams to a hardware-backed checker before deployment.